% Отчет по лабораторной работе №7
% Дисциплина: Системы управления базами данных
% Оформление выполнено по шаблону example.pdf
\documentclass[a4paper,12pt]{article}

\usepackage[T2A]{fontenc}
\usepackage[utf8]{inputenc}
\usepackage[russian]{babel}
\usepackage{tempora}
\usepackage[a4paper,margin=2cm]{geometry}
\usepackage{graphicx}

% Размер основного текста как в шаблоне: 14 pt, интервал 18.4 pt
\renewcommand{\normalsize}{\fontsize{14}{18.4}\selectfont}
\normalsize
\renewcommand{\large}{\fontsize{16.8}{20.16}\selectfont}
\renewcommand{\small}{\fontsize{12}{15}\selectfont}

\setcounter{secnumdepth}{2}

% Как в шаблоне: абзацы без отступа первой строки, без доп. интервала между абзацами
\setlength{\parindent}{0pt}
\setlength{\parskip}{0pt}

% Заголовки разделов: полужирные, того же кегля, что и текст; нумерация "1.", "1.1."
\makeatletter
\renewcommand\section{\@startsection{section}{\z@}{\z@}%
  {14pt \@plus 4pt \@minus 4pt}{7pt}{\normalsize\bfseries}}
\renewcommand\subsection{\@startsection{subsection}{\z@}{\z@}%
  {14pt \@plus 4pt \@minus 4pt}{7pt}{\normalsize\bfseries}}
\makeatother

% Изображение с подписью "Рисунок N -- ..."
\newcounter{fig}
\newcommand{\pic}[2]{%
  \begin{center}
    \stepcounter{fig}%
    \includegraphics[width=0.95\textwidth]{#1}\\[6pt]
    Рисунок~\arabic{fig}~-- #2
  \end{center}
}

% Жирная метка с новой строки (процедуры, функции, пункты проверки)
\newcommand{\grouplabel}[1]{\par\medskip\noindent\textbf{#1}\par\smallskip}

\begin{document}

% ============================ ТИТУЛЬНЫЙ ЛИСТ ============================
\begin{titlepage}
\centering
\vspace*{0.28\textheight}
{\large Системы управления базами данных\par}
\vspace{8pt}
{\large Лабораторная работа №7\par}
\vfill
\begin{flushright}
Выполнил: Сахань В.С.\\
ФИТ-242
\end{flushright}
\vspace*{0.05\textheight}
\end{titlepage}

% ============================ 1. ТЕМА И ЦЕЛЬ ============================
\section{Тема и цель работы}

\textbf{Тема:} Разработка хранимых процедур и функций в СУБД PostgreSQL на языках SQL и PL/pgSQL.

\textbf{Цель работы:} Приобретение практических навыков проектирования, реализации и тестирования серверных подпрограмм (хранимых процедур и пользовательских функций) для обеспечения целостности данных и инкапсуляции бизнес-логики предметной области «Кинотеатр».

\textbf{Индивидуальный вариант:} Вариант 17 («Кинотеатр»).

% ============================ 2. ФОРМУЛИРОВКИ ============================
\section{Формулировка процедур и функций}

В рамках работы разработаны три хранимые процедуры и две пользовательские функции предметной области «Кинотеатр».

\subsection{Хранимые процедуры}

\grouplabel{sp\_sell\_ticket (Оформление продажи билета)}

\textbf{Назначение:} Оформляет покупку билета кассиром на указанный сеанс. Автоматически определяет базовую цену сеанса, применяет скидку клиента (если билет оформляется по предзаказу), регистрирует билет в таблице \texttt{Tickets} и переводит связанную бронь в таблице \texttt{Preorders} в статус «выкуплен».

\textbf{Язык реализации:} plpgsql.

\grouplabel{sp\_cancel\_preorder (Аннулирование бронирования)}

\textbf{Назначение:} Реализует регламент отмены брони клиентом или при истечении срока выкупа. Переводит статус предзаказа в таблице \texttt{Preorders} в «отменен» и обнуляет привязку к брони (\texttt{id\_preorder = NULL}) у зарезервированных билетов в таблице \texttt{Tickets}.

\textbf{Язык реализации:} plpgsql.

\grouplabel{sp\_apply\_session\_discount (Пакетное изменение стоимости сеанса)}

\textbf{Назначение:} Корректирует базовую стоимость сеанса в таблице \texttt{Sessions} на заданный процент скидки (например, при утренних сеансах или акциях).

\textbf{Язык реализации:} sql.

\subsection{Хранимые функции}

\grouplabel{fn\_get\_session\_revenue (Расчет фактической выручки сеанса)}

\textbf{Назначение:} Скалярная функция, вычисляющая суммарный финансовый доход от проданных билетов на конкретный сеанс по его идентификатору \texttt{id\_session}.

\textbf{Возвращаемое значение:} NUMERIC(10, 2) (если билеты еще не проданы, возвращает 0.00).

\textbf{Язык реализации:} plpgsql.

\grouplabel{fn\_get\_movie\_sessions\_status (Анализ занятости зала)}

\textbf{Назначение:} Табличная функция (\texttt{RETURNS TABLE}), формирующая сводную ведомость по сеансам выбранного фильма: дата проведения, зал, общая вместимость, количество купленных билетов и число свободных мест для оперативной работы кассы.

\textbf{Возвращаемое значение:} Таблица со списком параметров сеансов (\texttt{session\_id}, \texttt{session\_time}, \texttt{hall\_name}, \texttt{total\_capacity}, \texttt{sold\_tickets}, \texttt{available\_seats}).

\textbf{Язык реализации:} sql.


% ============================ 3. СКРИПТ СОЗДАНИЯ ============================
\section{SQL-скрипт создания процедур и функций}

Серверные подпрограммы создаются командами \texttt{CREATE OR REPLACE PROCEDURE} и \texttt{CREATE OR REPLACE FUNCTION}. Скрипты сохранены в каталоге \texttt{Scripts} и выполнены в СУБД PostgreSQL.

\grouplabel{1. Процедура оформления продажи билета (\texttt{sp\_sell\_ticket})}

\pic{Screenshots/Code01.png}{Код создания процедуры продажи билета}
\pic{Screenshots/CodeResult01.png}{Результат работы кода создания процедуры продажи билета}

\grouplabel{2. Процедура аннулирования бронирования (\texttt{sp\_cancel\_preorder})}

\pic{Screenshots/Code02.png}{Код создания процедуры аннулирования бронирования}
\pic{Screenshots/CodeResult02.png}{Результат работы кода создания процедуры аннулирования бронирования}

\grouplabel{3. Процедура пакетного изменения стоимости сеанса (\texttt{sp\_apply\_session\_discount})}

\pic{Screenshots/Code03.png}{Код создания процедуры изменения стоимости сеанса}
\pic{Screenshots/CodeResult03.png}{Результат работы кода создания процедуры изменения стоимости сеанса}

\grouplabel{4. Функция расчета выручки сеанса (\texttt{fn\_get\_session\_revenue})}

\pic{Screenshots/Code04.png}{Код создания функции расчета выручки сеанса}
\pic{Screenshots/CodeResult04.png}{Результат работы кода создания функции расчета выручки сеанса}

\grouplabel{5. Табличная функция анализа занятости зала (\texttt{fn\_get\_movie\_sessions\_status})}

\pic{Screenshots/Code05.png}{Код создания табличной функции анализа занятости зала}
\pic{Screenshots/CodeResult05.png}{Результат работы кода создания табличной функции анализа занятости зала}

% ============================ 4. ТЕСТИРОВАНИЕ ============================
\section{Тестирование процедур и функций}

Каждая подпрограмма проверена отдельным тестовым скриптом: создается тестовая запись, вызывается процедура (\texttt{CALL}) или функция, после чего выборкой контролируется итоговое состояние данных.

\grouplabel{1. Проверка процедуры продажи билета}

Создается предзаказ, затем вызывается \texttt{sp\_sell\_ticket} с размером скидки 10\,\% и выборкой проверяется запись о билете и статус брони.

\pic{Screenshots/Code06.png}{Код проверки процедуры продажи билета}
\pic{Screenshots/CodeResult06.png}{Результат работы кода проверки процедуры продажи билета}

\grouplabel{2. Проверка процедуры аннулирования бронирования}

Создается тестовый предзаказ, вызывается \texttt{sp\_cancel\_preorder}, после чего проверяется статус записи в таблице \texttt{Preorders}.

\pic{Screenshots/Code07.png}{Код проверки процедуры аннулирования бронирования}
\pic{Screenshots/CodeResult07.png}{Результат работы кода проверки процедуры аннулирования бронирования}

\grouplabel{3. Проверка процедуры изменения стоимости сеанса}

Фиксируется стоимость сеанса до вызова \texttt{sp\_apply\_session\_discount} (скидка 20\,\%) и после него.

\pic{Screenshots/Code08.png}{Код проверки процедуры изменения стоимости сеанса}
\pic{Screenshots/CodeResult08.png}{Результат работы кода проверки процедуры изменения стоимости сеанса}

\grouplabel{4. Проверка функции расчета выручки сеанса}

С помощью \texttt{fn\_get\_session\_revenue} вычисляется фактическая выручка по сеансу с учетом проданных билетов.

\pic{Screenshots/Code09.png}{Код проверки функции расчета выручки сеанса}
\pic{Screenshots/CodeResult09.png}{Результат работы кода проверки функции расчета выручки сеанса}

\grouplabel{5. Проверка табличной функции анализа занятости зала}

Функция \texttt{fn\_get\_movie\_sessions\_status} возвращает ведомость по сеансам фильма: вместимость зала, проданные билеты и свободные места.

\pic{Screenshots/Code10.png}{Код проверки табличной функции анализа занятости зала}
\pic{Screenshots/CodeResult10.png}{Результат работы кода проверки табличной функции анализа занятости зала}

В ходе тестирования подтверждено, что все процедуры и функции работают в соответствии с назначением: продажа и отмена брони корректно изменяют состояние данных, скидки применяются к ценам, выручка и занятость зала рассчитываются без ошибок.

\end{document}
